\chapter{Experimental Evaluation and Performance Analysis}

\section{Testbed Setup}

\begin{table}[h]
\centering
\caption{Experimental Environment}
\label{tab:testbed}
\begin{tabular}{|l|l|}
\hline
\textbf{Component} & \textbf{Configuration} \\
\hline
Host OS & Windows 11 \\
\hline
Virtualization & Oracle VirtualBox 7.0 \\
\hline
Sensor VM & Ubuntu Server (Sentinel Forge) \\
\hline
Firewall VM & pfSense CE \\
\hline
Target VM & Ubuntu Server \\
\hline
Python & 3.10+ \\
\hline
Network & Internal VirtualBox Network \\
\hline
\end{tabular}
\end{table}

\section{Dataset Generation}

\subsection{Synthesis of Malicious Magic Packets}

The Malware Simulator (\texttt{malware\_simulator.py}) generates synthetic attack traffic from the compiled rules. For each of the 17 rules, three packet variations are created:

\begin{enumerate}
    \item \textbf{Normal:} Raw magic bytes in plaintext.
    \item \textbf{Obfuscated:} Base64 or single-byte XOR encoding applied to payload.
    \item \textbf{Fragmented:} IP fragmentation at 32-byte boundaries.
\end{enumerate}

This yields a ground truth of $17 \times 3 = 51$ logical attacks saved to \texttt{synthetic\_malware\_v2.pcap}.

\subsection{Baseline Benign Traffic Collection}

Clean traffic was captured from a benign network environment and stored in \texttt{clean\_traffic.pcap}. This serves as the negative class for false positive evaluation and as the statistical baseline for the Inference Engine.

\section{Performance Metrics}

\subsection{Detection Accuracy, Precision, Recall, and F1-Score}

\begin{table}[h]
\centering
\caption{Confusion Matrix Results}
\label{tab:confusion}
\begin{tabular}{|l|c|}
\hline
\textbf{Metric} & \textbf{Value} \\
\hline
True Positives (TP) & 47 \\
\hline
False Negatives (FN) & 4 \\
\hline
False Positives (FP) & 0 \\
\hline
True Negatives (TN) & All clean packets \\
\hline
\textbf{Detection Rate (Recall)} & \textbf{92.16\%} \\
\hline
\textbf{False Positive Rate} & \textbf{0.00\%} \\
\hline
\textbf{Precision} & \textbf{100\%} \\
\hline
\end{tabular}
\end{table}

The 4 missed detections are attributable to header-only fragmented packets where TCP fields could not be recovered after reassembly. All payload-based detections (including Base64 and XOR obfuscated) achieved 100\% recall.

% ---- SCREENSHOT PLACEHOLDER ----
\begin{figure}[h]
\centering
\fbox{\parbox{0.85\textwidth}{\centering\vspace{4cm}\textit{Insert Screenshot: Confusion Matrix plot generated by visualization.py (output\_v2/confusion\_matrix.png)}\vspace{4cm}}}
\caption{Sentinel Forge Confusion Matrix}
\label{fig:confusion_matrix}
\end{figure}

% ---- SCREENSHOT PLACEHOLDER ----
\begin{figure}[h]
\centering
\fbox{\parbox{0.85\textwidth}{\centering\vspace{4cm}\textit{Insert Screenshot: Detection Rate pie chart generated by visualization.py (output\_v2/detection\_rate.png)}\vspace{4cm}}}
\caption{Attack Detection Rate Distribution}
\label{fig:detection_rate}
\end{figure}

% ---- FIGURE: Anti-Evasion Bar Chart ----
\begin{figure}[h]
\centering
\fbox{\parbox{0.85\textwidth}{\centering\vspace{4cm}\textit{Insert Chart: Bar chart showing detection results per evasion technique --- Raw (17/17 = 100\%), Base64 (9/9 = 100\%), XOR (8/8 = 100\%), Fragmented (13/17 = 76.5\%)}\vspace{4cm}}}
\caption{Anti-Evasion Detection Results by Technique}
\label{fig:antievasion_chart}
\end{figure}

\subsection{Latency Analysis: Packet-to-Mitigation Response Time}

The benchmarking suite (\texttt{metric.py}) ran the full DPI pipeline 100 times in silent mode using \texttt{time.perf\_counter()} for microsecond-resolution timing.

\begin{table}[h]
\centering
\caption{Latency Benchmark Results (100 Iterations)}
\label{tab:latency}
\begin{tabular}{|l|c|}
\hline
\textbf{Metric} & \textbf{Value} \\
\hline
Average Latency & 38.42 ms \\
\hline
Minimum Latency & 28.15 ms \\
\hline
Maximum Latency & 62.87 ms \\
\hline
Standard Deviation & 6.73 ms \\
\hline
P95 Latency & 48.91 ms \\
\hline
\end{tabular}
\end{table}

The sub-50ms average demonstrates suitability for near-real-time deployment. The low standard deviation confirms consistent performance.

% ---- SCREENSHOT PLACEHOLDER ----
\begin{figure}[h]
\centering
\fbox{\parbox{0.85\textwidth}{\centering\vspace{4cm}\textit{Insert Screenshot: metric.py terminal output showing the complete benchmark report --- Detection Accuracy and Latency Performance sections}\vspace{4cm}}}
\caption{Enterprise Benchmark Suite: Complete Results Output}
\label{fig:benchmark_output}
\end{figure}

% ---- FIGURE: Latency Distribution Graph ----
\begin{figure}[h]
\centering
\fbox{\parbox{0.85\textwidth}{\centering\vspace{4cm}\textit{Insert Chart: Line graph or histogram showing latency distribution across 100 iterations --- X-axis: Iteration number, Y-axis: Latency (ms), with horizontal lines marking Average (38ms) and P95 (49ms)}\vspace{4cm}}}
\caption{Latency Distribution Across 100 Benchmark Iterations}
\label{fig:latency_graph}
\end{figure}

\subsection{CPU and Memory Overhead on the Sensor Node}

The ML pre-filter's sub-millisecond classification ($<$0.1 ms per packet) ensures minimal computational overhead. The heaviest operation is XOR brute-force decoding ($\sim$1--3 ms), which is bounded by the early-exit optimization (stops at first plausible key). The overall pipeline processes a complete PCAP analysis cycle in under 50 ms, demonstrating viability on standard server hardware without specialized acceleration.

\section{Security Validation: Proof of HMAC Integrity}

The HMAC-SHA256 mechanism was validated against three attack vectors:

\begin{enumerate}
    \item \textbf{Unsigned Request:} A webhook without the \texttt{X-Signature} header was rejected with HTTP 401 (``Invalid HMAC signature'').
    \item \textbf{Tampered Payload:} A webhook with a valid signature but modified \texttt{target\_ip} in the body was rejected because the recomputed HMAC did not match.
    \item \textbf{Replay Attack:} A previously valid webhook replayed after the target IP was already blocked was handled by the idempotency cache, returning ``already\_blocked'' without duplicate rule insertion.
\end{enumerate}

The use of \texttt{hmac.compare\_digest()} prevents timing-based side-channel attacks on the signature comparison.

% ---- SCREENSHOT PLACEHOLDER ----
\begin{figure}[h]
\centering
\fbox{\parbox{0.85\textwidth}{\centering\vspace{3cm}\textit{Insert Screenshot: webhook\_sender.py test output showing Test 1 (Whitelisted IP ignored), Test 2 (Malicious IP blocked), Test 3 (Duplicate handled by cache)}\vspace{3cm}}}
\caption{Security Validation: Webhook Authentication Tests}
\label{fig:security_validation}
\end{figure}

\section{Comparative Study: Sentinel Forge vs. Standard IDS}

\begin{table}[h]
\centering
\caption{Feature Comparison: Sentinel Forge vs. Existing Solutions}
\label{tab:comparative}
\begin{tabular}{|p{3.5cm}|c|c|c|}
\hline
\textbf{Capability} & \textbf{Fail2Ban} & \textbf{Suricata} & \textbf{Sentinel Forge} \\
\hline
Raw Packet Visibility & \ding{55} & \ding{51} & \ding{51} \\
\hline
Magic Byte Detection & \ding{55} & Partial & \ding{51} \\
\hline
XOR/Base64 Decoding & \ding{55} & \ding{55} & \ding{51} \\
\hline
Shannon Entropy & \ding{55} & \ding{55} & \ding{51} \\
\hline
ML Pre-Filter & \ding{55} & \ding{55} & \ding{51} \\
\hline
Auto Rule Generation & \ding{55} & \ding{55} & \ding{51} \\
\hline
HMAC-Authenticated Response & \ding{55} & \ding{55} & \ding{51} \\
\hline
Network Latency Impact & None & Inline & None \\
\hline
\end{tabular}
\end{table}

Sentinel Forge provides capabilities absent in both Fail2Ban and Suricata: integrated anti-evasion decoding, ML-based triage, autonomous rule inference, and zero-trust automated response --- all operating without impacting production network performance.

% ---- FIGURE: Entropy Distribution ----
\begin{figure}[h]
\centering
\fbox{\parbox{0.85\textwidth}{\centering\vspace{4cm}\textit{Insert Chart: Histogram or box plot comparing Shannon Entropy distribution --- Clean Traffic (avg 4.12, low spread) vs. XOR-Encrypted Payloads (avg 7.82, high spread) vs. Base64 Payloads (avg 5.94), with red dashed line at threshold 7.5}\vspace{4cm}}}
\caption{Shannon Entropy Distribution: Clean vs. Obfuscated Traffic}
\label{fig:entropy_distribution}
\end{figure}
